\subsection{A stronger local theorem from report 105: defer the lower forest}
\label{marked-cut:deferred-review}

Report~105 arrived while the preceding implementation was being validated.
Its delivered code remains a separate research package. A subsequent review
checks the following stronger local argument and replays its positive
packets with the current native independent verifier. This section records
that argument and the integration boundary; the maintained default is still
the source-first maximum-passage retry of the preceding section.

\paragraph{The source condition.}
Let $G=(V,E)$ be the zero/one-handle multigraph and let $N$ be the full
nonnegative occurrence matrix before a lower forest is removed. Its columns
are $E\sqcup J$, where $J$ contains auxiliary bundle columns, if present.
For each row, assume its \emph{uncapped} restriction to $E$ is a mod-two
cycle of $G$. A literal closed attaching walk has this property. Loops
contribute twice at their vertex. A row may restrict to the empty walk.
Saturation at two is legitimate for singleton tests but cannot be used to
establish this parity hypothesis: a count of three has odd parity.

\begin{proposition}[Report 105: peeling leaves a spanning forest]
Run column-leaf peeling on $N$, possibly stopping before all rows are
removed. Let $P\subseteq E$ be the graph columns used as pivots. Then
$G-P$ has exactly the original connected components of $G$.
\end{proposition}
\begin{proof}
Induct over the pivot order. An auxiliary pivot removes no graph edge.
For a graph pivot $e$ in row $r$, all earlier pivot columns have zero entry
in $r$: they were column leaves when their different rows were removed.
Thus $r$ is still a mod-two cycle in the graph with earlier pivots deleted.
It crosses $e$ once. A bridge separates a component into a cut, and every
mod-two cycle crosses that cut evenly, so $e$ cannot be a bridge of the
current graph. Removing it preserves the component partition. Induction
proves the assertion and hence the existence of a spanning forest avoiding
all graph pivots.
\end{proof}

\begin{proposition}[Report 105: complete lower-forest search is unnecessary]
Under this cycle hypothesis, full column-leaf peeling succeeds if and only
if some spanning lower forest can be removed so that complete raw positive
singleton cancellation succeeds. The construction requires one full
peeling pass and one spanning-forest pass, without enumerating forests.
\end{proposition}
\begin{proof}
If full peeling succeeds, choose a spanning forest in $G-P$ by the previous
proposition. Its edges are not pivots. Discarding these extra columns
changes none of the pivot-column sums, so the same complete peeling trace
works after contracting this forest.

Conversely, complete raw positive cancellation after any forest deletion
implies complete column peeling by the stopping-set theorem already proved
for the nonnegative cancellation model. Restore the forest columns. The
same pivots still have the same live column sums, so their sequence remains
valid in the full matrix. The order independence of the terminal peeling
residual then gives success of the full peeling process.
\end{proof}

This changes discovery order, not geometric execution order. The emitted
certificate first contracts the discovered forest and then performs the
usual source-relative one/two cancellations. It assumes no commutation of
geometric moves. A nonempty terminal full-matrix stopping set obstructs
complete raw cancellation for \emph{every} lower forest, since deleting
columns cannot destroy its zero-or-at-least-two column sums. It remains an
obstruction only to this fixed raw model, not to being a handlebody.

The same argument on a partial peeling trace produces a forest attaining
the smallest terminal column-peeling residual over all lower forests.
It does not maximize the length of arbitrary positive-fill partial traces:
the all-ones two-by-two matrix has no column leaf but allows one singleton
pivot before the remaining entry becomes two. The exact conclusion is
complete-success equivalence and optimality of the peeling residual.

\paragraph{Fresh review evidence and a specific next integration.}
An isolated copy of report~105 over its pinned
\texttt{c6d0ec8aa0b6} baseline passes all eleven delivered deferred-forest
and boundary tests in 8.864 seconds. These include 120 seeded closed-walk
instances compared against exhaustive forest choices. The initial local
attempt lacked Regina on its import path; the complete rerun includes the
required package. No delivered file was changed.

Separately, all nine supplied whole-cut packets and eight supplied core
packets pass the current independent native checkers. The core replay also
reconstructs their source geometry and the original queue top reduction.
Eight of those whole-cut sources already complete under the newly
integrated adaptive policy. The remaining source,
\path{pachner23_5_8_2/mixed}, has a valid delivered proof and agrees with
our fresh literal census, but still receives an inconclusive result from
the maintained producer. That packet is an integration target, not a
111th production success in the preceding audit.

A new native research probe uses the current maximum-weight top reduction
on all 120 sources, checks every retained raw row's uncapped vertex parity,
peels the full core matrix, and selects a forest avoiding its pivots. It
passes that forest into the ordinary native core API and independently
replays every success. The result is 118 complete bare cores, six more than
canonical lower forests on the same top reductions. The two remaining
stopping sets occur on \path{nested_1_4_6/interior_links} and
\path{interior_loop/interior_links}. These are bare-core counts; added
annular rows can invalidate a core-only pivot. A whole-cut integration
must assemble the complete relevant boundary matrix before deferred
selection and retain source authentication, top-reduction semantics,
omitted-column lifting and shared query budgets.

\paragraph{The distinct exceptional-interface result.}
Report~105 also extends report~99's two-row column selector to a finite
parameterized problem. Let $u$ count rows meeting an active column with
support on more than two rows, where active means that some entry is one.
The delivered reduction uses at most $\operatorname{Bell}(u+1)$ polynomial
independent-representatives calls: two-row implications are compressed,
partial partitions of the exceptional interface specify seed signatures,
and graphic/partition matroid intersection chooses compatible tree edges.
The partial partitions correspond to ordinary partitions after adjoining
one distinguished element, and
$\operatorname{Bell}(u+1)\leq(u+1)^u$. This yields
$2^{O(u\log(u+2))}\operatorname{poly}(N)$ for its finite graph/matrix model.
The original native triple-support specialization bounds $u\leq3k$ for
$k$ such columns; added annular rows require their actual supports.

An isolated rerun of its independent exhaustive oracle checks 50,658
finite inputs, including 46,656 exhaustive small inputs and 4,000 seeded
cases. It reproduces 5,183 successful independently replayed tree/peeling
certificates with zero discrepancies. This finite evidence is not a source
adapter, a native parameter bound or an independent formalization of the
whole manuscript. The coupled geometric adapter and a global bound on
$u$ remain future work. The deferred theorem and this parameterized
selector remove local search variables under explicit hypotheses; neither
establishes general quasi-polynomial unknot recognition.
